\documentclass[10pt,a4paper]{amsart}
\usepackage[margin=19mm]{geometry}
\usepackage{amsmath,amssymb,mathtools}
\usepackage[scr=rsfs,cal=euler]{mathalfa}
\usepackage{xfrac}
\usepackage[unicode,hidelinks,bookmarks=false]{hyperref}
\newcommand{\ie}{i.e.\ }
\newcommand{\cf}{cf.\ }
\newcommand{\resp}{respectively\ }
\newcommand{\Subsection}{\subsection}
\providecommand{\nobreakdash}{\nobreak\hbox{-}}
\setlength{\emergencystretch}{2em}
\multlinegap0pt
\usepackage{amssymb}
\usepackage{tikz}
\newtheorem{proposition}{Proposition}
\newcommand{\N}{\mathbb{N}}
\newcommand{\Z}{\mathbb{Z}}
\title{Parametrized complexity of relations between multidimensional subshifts}
\author{Nicanor Carrasco-Vargas, Benjamin Hellouin de Menibus,, R\'emi Pallen}
\date{Source: arXiv:2509.24343v2 (CC BY 4.0)}
\begin{document}
\maketitle

\noindent\emph{Source and licence.} Parametrized complexity of relations between multidimensional subshifts. Version arXiv:2509.24343v2 (CC BY 4.0), distributed by arXiv under the Creative Commons Attribution 4.0 International licence, http://creativecommons.org/licenses/by/4.0/, as recorded in the arXiv licence metadata for this exact version and shown on the abstract page https://arxiv.org/abs/2509.24343v2. Full licence notice: \texttt{licenses/arxiv-2509.24343-CC-BY-4.0.txt}.\par\smallskip

\noindent\textit{Editorial scope.} Counted unit: the article's proposition hardness-conjugacy-with-effective-input (source0.tex lines 706--708). The article's complete original proof of it is delivered with it (source0.tex lines 709--718, one \texttt{proof} environment of 10 lines). No mathematics is written, repaired or re-derived: the statement and the proof are the article's own lines, in the article's own notation, and no retained line is substituted or re-flowed.  No further source-local context is needed: every cross-reference of every retained line resolves inside the two delivered spans.  Nothing was omitted from any retained span, so the package carries no omission record.  No retained line contains a citation, so the wrapper carries no bibliography and no bibliography entry.  No retained line contains a figure environment, an image-inclusion command or drawing code, so no image is delivered and the package declares no asset dependency.  Editorial formatting only, in addition: an AMS \texttt{amsart} preamble that restates the article's own declarations and macros used by the retained lines, namely the environments \texttt{proposition} on separate counters, together with the standard packages those lines need.  The retained environments print in this document as Proposition 1 in order of appearance, whereas the article prints the same items with the numbering of its own full document.  The complete native source of the article as distributed in the arXiv e-print for this exact version is bundled unchanged as \texttt{sources/185941/source0.tex}.
\begin{proposition}\label{hardness-conjugacy-with-effective-input}
    Let $Y$ be a nonempty effective subshift. Then the conjugacy problem $X\simeq Y$ with effective inputs is $D(\Sigma_1^0)$-hard. 
\end{proposition}

\begin{proof}
Let $A\subset\N$ be a $D(\Sigma_1^0)$-complete set, and $B$ and $C$ be $\Sigma_1^0$-sets such that $C\subseteq B$ and $A=B\smallsetminus C$. 
For each $n\in\N$, we define $X_n$ on alphabet $\{0,1\}$ as follows: if $n\in B$, we forbid the symbol $0$, and if $n\in C$, we forbid the symbol $1$. This set of forbidden patterns is recursively enumerable from $n$, so $(X_n)_{n\in \N}$ is a computable sequence of effective subshifts.
\[X_n=\begin{cases}
\{1\}^{\Z^d} &\ n\in B, \ n\notin C\\
\emptyset &\ n \in B, \  n\in C \\
\{0,1\}^{\mathbb{Z}^d} &\ n\notin B, \ n\notin C
\end{cases}\]
Notice that $Y\simeq Y\times X_n$ exactly when $|X_n|=1$ (as $Y\times X_n$ is otherwise empty or has too much entropy). This is equivalent to $n\in A$, which is a $D(\Sigma_1^0)$-complete problem. 
\end{proof}

\end{document}
